\section{Introduction: Exoplanet Vetting and the Signal Ambiguity Challenge} \label{sec:introduction}

\subsection{Context of Modern Wide-Field Space Photometry} \label{sec:intro_context}
The search for transiting exoplanets has transitioned from target-specific ground-based observations to wide-field space-based photometric surveys. Space missions, such as the CoRoT satellite \citep{Baglin2006}, the Kepler Space Telescope \citep{Borucki2010, Koch2010, Batalha2013}, the K2 mission \citep{Howell2014}, the Transiting Exoplanet Survey Satellite \citep[TESS;][]{Ricker2015, Vanderspek2018, Guerrero2021}, the CHEOPS mission \citep{Benz2021}, and the upcoming PLATO mission \citep{Rauer2014}, monitor hundreds of thousands to millions of stars simultaneously. By measuring stellar flux continuously at high cadence (e.g., 2-minute and 20-second integrations for TESS \citealt{Jenkins2016}; 30-minute and 1-minute integrations for Kepler), these missions generate vast quantities of time-series photometric data that have revolutionized our understanding of exoplanetary demographics and system architectures \citep{Howard2012, Petigura2013, Winn2015, Fulton2017}.

To detect planetary transits---periodic brightness dips caused by a planet crossing the disk of its host star---automated search algorithms are applied to detrended light curves. Standard search algorithms include the Box-fitting Least Squares \citep[BLS;][]{Kovacs2002}, Transit Least Squares \citep[TLS;][]{Hippke2019}, Quasi-periodic Automated Transit Search \citep[QATS;][]{Carter2013}, and optimized detection efficiency variants \citep{Tingley2005, Ofir2014}. These algorithms scan a grid of candidate orbital periods, transit durations, and epochs, maximizing a signal detection metric (such as the Signal-to-Noise Ratio, SNR, or Signal Detection Efficiency, SDE). Signals exceeding a predefined detection threshold (originating with the Kepler pipeline at $\text{SNR} \ge 7.1$ \citealt{Jenkins2002, Jenkins2010}; or with TLS at $\text{SDE} \ge 6.0$ \citealt{Hippke2019}) are flagged as Threshold Crossing Events (TCEs) or candidate events.

However, only a small fraction of TCEs correspond to true planetary systems. In wide-field space photometry, false positive rates often exceed 90\% \citep{Fressin2013}, as evaluated in the final Kepler catalog releases \citep{Mullally2015, Coughlin2016, Thompson2018}. These false positives arise from three principal mechanisms:
\begin{enumerate}
\item \textbf{Instrumental Anomalies}: Pointing jitter, thermal drifts, rolling band noise, solar wind interactions, pixel-level charge leaks, and Presearch Data Conditioning (PDC) artifacts \citep{Smith2012, Stumpe2012}.
\item \textbf{Astrophysical False Positives}: Eclipsing binary (EB) stars \citep{Slawson2011, Prsa2011, Kirk2016} or blended background eclipsing binaries (BEBs) blended within the photometric aperture \citep{Santerne2012, Bryson2013, Sullivan2015, Colon2012}. Because TESS has large pixels ($21''$ per pixel) and TESS Input Catalog (TIC) targets span crowded galactic fields \citep{Stassun2018, Stassun2019}, flux contamination from nearby stars frequently blends eclipsing binary signals into the target star's aperture.
\item \textbf{Stellar Variability and Noise}: Intrinsic stellar pulsations, convective granulation noise, and magnetic activity \citep{Meunier2010, McQuillan2013, McQuillan2014, Reinhold2013, Aigrain2012}.
\end{enumerate}

Due to the sheer volume of candidates (e.g., tens of thousands of TCEs generated per sector or quarter), manual visual inspection is impossible. Consequently, automated vetting systems are required to filter out false positives and rank candidates before allocating limited ground-based follow-up resources (such as high-precision radial velocity spectrographs or adaptive optics imaging).

\subsection{Automated Vetting Evolution: History and Context} \label{sec:vetting_history}
The historical evolution of automated exoplanet transit vetting has progressed through three distinct eras, moving from heuristic, rule-based expert systems to complex machine learning classifiers, and finally to computationally intensive Bayesian simulations (as shown in Figure~\ref{fig:evolution}).

\begin{figure*}[htbp]
\centering
\begin{tikzpicture}[node distance=2.2cm, auto, >=latex', thick,
    block/.style={rectangle, draw, fill=blue!10, text width=3.3cm, text centered, rounded corners, minimum height=1.3cm},
    line/.style={draw, ->, shorten >=1pt}]
    
    \node [block] (robo) {
        \textbf{Kepler Robovetter}\\
        \small Heuristic Rules\\
        \small Low Complexity
    };
    \node [block, right of=robo, node distance=4.5cm] (ml) {
        \textbf{Astronet / Autovetter}\\
        \small High-Capacity ML\\
        \small Overfitting Susceptible
    };
    \node [block, right of=ml, node distance=4.5cm] (bayes) {
        \textbf{Vespa / Triceratops}\\
        \small Bayesian MCMC FPP\\
        \small High CPU-Hours
    };
    \node [block, right of=bayes, node distance=4.5cm, fill=green!10] (tars) {
        \textbf{TARS v1 (RAI)}\\
        \small Parsimonious Epistemic\\
        \small Fast \& Interpretable
    };
    
    \path [line] (robo) -- (ml);
    \path [line] (ml) -- (bayes);
    \path [line] (bayes) -- (tars);
\end{tikzpicture}
\caption{Historical evolution of automated exoplanet vetting methodologies, showing the progression from heuristic rules to modern parsimonious indexing.}
\label{fig:evolution}
\end{figure*}

\subsubsection{Rule-Based Heuristics}
The early Kepler pipelines utilized the Robovetter \citep{Coughlin2016}. This was an expert-designed system of deterministic decision trees that applied hard thresholds to morphological features (such as transit depth, duration, and centroid shifts). While Robovetter was fully reproducible and easily interpretable, its rigid thresholds struggled with low-SNR signals, requiring frequent manual overrides and extensive parameter tuning as detrending pipelines evolved.

\subsubsection{High-Capacity Machine Learning Ensembles}
To handle near-threshold detections, the community developed high-capacity classifiers. Autovetter \citep{McCauliff2015} introduced random forests to exoplanet vetting, utilizing dozens of engineered features. This was followed by deep learning approaches such as Astronet \citep{Shallue2018}, which employed 1D convolutional neural networks (CNNs) to classify phase-folded light curves directly. Subsequent architectures extended deep learning vetting to domain adaptation \citep{Ansdell2018}, multi-sector TESS classification \citep{Yu2019, Osborn2020, Olmschenk2021}, K2 discovery \citep{Dattilo2019}, and ensemble planet validation \citep{Armstrong2021}. While these classifiers improved ranking accuracy, they brought significant challenges:
\begin{itemize}
\item \textbf{Overfitting and Redundancy}: Ensembles typically utilize dozens of highly correlated features, making them susceptible to overfitting under dataset shift.
\item \textbf{Uncalibrated Output Probabilities}: High-capacity models frequently yield uncalibrated confidence scores \citep{Guo2017, Niculescu2005}, making it difficult for observers to interpret the outputs as true physical probabilities.
\item \textbf{Lack of Interpretability}: Deep neural networks behave as black boxes, providing no physical justification for candidate rejection.
\end{itemize}

\subsubsection{Bayesian Stellar Population Simulations}
To obtain physically meaningful probability estimates, Bayesian tools such as Vespa \citep{Morton2012, Morton2015, Morton2016}, PASTIS \citep{Diaz2014, Santerne2015}, BLENDER \citep{Torres2011}, and Triceratops \citep{Giacalone2021} calculate the False Positive Probability (FPP) of a candidate. They simulate millions of synthetic stellar populations (including blended background stars and orbiting binaries) to compute the relative likelihood of the data under different transit models. While highly rigorous, these tools are computationally expensive, requiring minutes to hours of CPU time per candidate. This creates a computational bottleneck when processing large-scale catalogs.

\subsection{The Transit Vetting Pipeline Flow and Bottlenecks} \label{sec:vetting_flow}
Automated vetting is not a single decision, but rather the culmination of a multi-stage data processing pipeline. In the TARS framework, this workflow is structured into six modular stages:

\begin{figure*}[htbp]
\centering
\begin{tikzpicture}[auto, >=latex', thick,
    block/.style={rectangle, draw, fill=blue!5, text width=4.5cm, text centered, rounded corners, minimum height=0.7cm},
    descr/.style={rectangle, text width=5.5cm, align=left, minimum height=0.7cm},
    line/.style={draw, ->, shorten >=1pt}]
    
    \node (raw) {\textbf{Raw Light Curve}};
    
    \node [block, below of=raw, node distance=1.1cm] (s1) {\textbf{Stage 1: Signal Conditioning}};
    \node [descr, right of=s1, node distance=5.5cm] (d1) {Detrending (splines/CBVs), systematic removal};
    
    \node [block, below of=s1, node distance=1.1cm] (s2) {\textbf{Stage 2: Transit Detection}};
    \node [descr, right of=s2, node distance=5.5cm] (d2) {Period searches (BLS/TLS), SNR thresholds};
    
    \node [block, below of=s2, node distance=1.1cm] (s3) {\textbf{Stage 3: Period Recovery}};
    \node [descr, right of=s3, node distance=5.5cm] (d3) {Multi-sector ephemeris fits, candidate families};
    
    \node [block, below of=s3, node distance=1.1cm] (s4) {\textbf{Stage 4: Morphology Vetting}};
    \node [descr, right of=s4, node distance=5.5cm] (d4) {Transit modeling, physical parameter checks};
    
    \node [block, below of=s4, node distance=1.1cm] (s5) {\textbf{Stage 5: Consistency Checks}};
    \node [descr, right of=s5, node distance=5.5cm] (d5) {Odd/even depths, centroid offsets};
    
    \node [block, below of=s5, node distance=1.1cm] (s6) {\textbf{Stage 6: Classification}};
    \node [descr, right of=s6, node distance=5.5cm] (d6) {Logistic model calibration, probability link};
    
    \path [line] (raw) -- (s1);
    \path [line] (s1) -- (s2);
    \path [line] (s2) -- (s3);
    \path [line] (s3) -- (s4);
    \path [line] (s4) -- (s5);
    \path [line] (s5) -- (s6);
\end{tikzpicture}
\caption{Modular layout of the 6-stage TARS automated exoplanet validation architecture, mapping data flow from raw flux inputs to final calibrated reliability probabilities.}
\label{fig:pipeline_flow}
\end{figure*}

The layout of the vetting architecture is summarized in Figure~\ref{fig:pipeline_flow}.
\begin{enumerate}
\item \textbf{Stage 1 (Signal Conditioning)}: Raw flux series are processed to remove long-term stellar rotational modulation and instrumental systematics \citep{Smith2012, Stumpe2012, Vanderburg2014, Luger2016, Luger2018, Aigrain2016}. This is achieved via high-pass filtering, iterative spline fitting, or projecting onto Co-trending Basis Vectors (CBVs). The conditioning stage must balance systematic removal against the preservation of long-period transits.
\item \textbf{Stage 2 (Transit Detection)}: The detrended light curve is searched for periodic transit signatures. Candidates are identified by finding the period $P$, epoch $T_0$, and duration $W$ that maximize the transit depth significance using algorithms such as BLS \citep{Kovacs2002} or TLS \citep{Hippke2019}.
\item \textbf{Stage 3 (Period Recovery \& Ephemeris Matching)}: Transit events from separate observation sectors or quarters are combined to fit a coherent linear ephemeris:
   \begin{equation}
   t_k = T_0 + k \cdot P
   \label{eq:ephemeris}
   \end{equation}
   This stage resolves the multi-peak periodogram structure \citep{Deeming1975, Roberts1987, Lomb1976, Scargle1982, VanderPlas2018, Dawson2010}, generating a list of candidate periods that match the observed transit timings.
\item \textbf{Stage 4 (Morphological Vetting)}: The light curve is folded at the candidate period and fitted with an analytical transit model \citep{Mandel2002, Gimenez2006, Kreidberg2015, Seager2003} to derive physical parameters, including the ratio of planet radius to stellar radius $R_p/R_*$, semi-major axis $a/R_*$, and impact parameter $b$.
\item \textbf{Stage 5 (Consistency Validation)}: Checks are executed to identify eclipsing binaries, such as comparing the depths of odd and even transits (which differ for eclipsing binaries with eccentric orbits or different secondary components) and measuring centroid offsets during transit to locate the true source of the dip.
\item \textbf{Stage 6 (Probabilistic Classification)}: Feature vectors containing metrics from all previous stages are fed into machine learning classifiers to predict the probability of the candidate being a true planet, evaluated via diagnostic metrics such as AUROC, AUPRC, DeLong statistics \citep{DeLong1988}, Youden's J index \citep{Youden1950}, and Matthews Correlation Coefficient \citep[MCC;][]{Matthews1975}.
\end{enumerate}

\subsection{Stellar Activity and Vetting Degeneracy} \label{sec:stellar_activity}
The primary astrophysical bottleneck in transit vetting is stellar activity. Active stars exhibit magnetic starspots on their photospheres \citep{Meunier2010, Aigrain2012}. As the star rotates, these spot groups cross the visible stellar hemisphere, causing quasi-periodic flux variations at the stellar rotation period $P_{\text{rot}}$ and its harmonics \citep{McQuillan2013, McQuillan2014, Reinhold2013}. Furthermore, starspot groups grow and decay on timescales of days to weeks, introducing amplitude and phase variations in the light curve.

This magnetic variability introduces two major challenges for transit detection:
\begin{enumerate}
\item \textbf{False Transits}: The spot groups and active regions can create temporary, localized flux dips that resemble transits when observed over a short baseline.
\item \textbf{Alias Frequency Clustering}: On active stars, stellar rotation creates quasi-periodic flux dips that can alias (fold) onto the true planetary transit period \citep{Dawson2010}. When performing transit searches, the algorithm cannot easily distinguish the true transit signal from stellar rotation harmonics ($P_{\text{rot}}/2, 2 P_{\text{rot}}$, etc.), producing multiple false candidate period peaks that survive standard vetting checks.
\end{enumerate}

When a transit search algorithm folds the light curve at these incorrect alias periods, the resulting phase-folded light curve can exhibit a high signal-to-noise ratio, confusing classical classifiers. The number of candidate periods that survive initial checks---often termed ``candidate branching'' or ``alias explosion''---increases exponentially, creating high signal recovery ambiguity.

\subsection{Aleatoric vs. Epistemic Uncertainty in Exoplanet Vetting} \label{sec:uncertainty_types}
To build robust classifiers, it is essential to distinguish between two types of uncertainty:
\begin{itemize}
\item \textbf{Aleatoric Uncertainty}: This represents the physical randomness inherent in the astronomical source. In vetting, it corresponds to whether the physical source is a planet, an eclipsing binary, or a blended background binary. This uncertainty is driven by physical degeneracies (e.g., a small star eclipsing a large star can produce the same transit depth as a planet transiting a Sun-like star) and is typically modeled using Bayesian false positive probability (FPP) packages like \texttt{vespa} \citep{Morton2012, Morton2015}.
\item \textbf{Epistemic Uncertainty}: This represents the uncertainty associated with the signal recovery process itself under noise and activity. It answers the question: \textit{How reliably can we recover the true period and ephemeris of the candidate given the observation window and the stellar activity profile?}
\end{itemize}

Many existing vetting pipelines primarily optimize classification performance rather than explicitly separating epistemic recovery ambiguity from astrophysical false-positive probability. They construct classifiers containing up to dozens of correlated features (such as ``family complexity'', which counts raw alias frequencies and periodogram peaks) to predict the class label $P(\text{Planet} \mid X)$. Because these classifiers are fit on historical catalogs, they are highly sensitive to dataset selection biases and cannot trace \textit{why} a candidate is flagged as ambiguous.

\subsection{Operationalizing a Multi-Dimensional Validation Protocol} \label{sec:contribution}

To address the limitations of conventional exoplanet vetting evaluations, this work operationalizes and evaluates a reproducible, multi-dimensional validation protocol for ambiguity-aware transit-vetting pipelines. Rather than evaluating a classifier solely on a single, local train/test split, our protocol establishes a multi-layered certification hierarchy (Figure~\ref{fig:validation_hierarchy}) designed to expose hidden operational limits, calibration decay under domain shifts, and statistical redundancies.

\begin{figure*}[htbp]
\centering
\begin{tikzpicture}[node distance=1.5cm, auto, >=latex', thick,
    block/.style={rectangle, draw, fill=blue!5, text width=6.5cm, text centered, rounded corners, minimum height=0.9cm},
    line/.style={draw, ->, shorten >=1pt}]
    
    \node [block] (acc) {\textbf{Layer 1: Discrimination \& Accuracy} \\ \small (Single-split AUROC \& F1 benchmarking)};
    \node [block, below of=acc, node distance=1.8cm] (cal) {\textbf{Layer 2: Probability Calibration} \\ \small (ECE and reliability curves by mission)};
    \node [block, below of=cal, node distance=1.8cm] (rob) {\textbf{Layer 3: Measurement Robustness} \\ \small (Stability under noise, jitter, and data loss)};
    \node [block, below of=rob, node distance=1.8cm] (dom) {\textbf{Layer 4: Domain Generalization} \\ \small (Cross-mission transfer: TESS $\rightarrow$ Kepler/K2)};
    \node [block, below of=dom, node distance=1.8cm] (nest) {\textbf{Layer 5: Nested Control Audits} \\ \small (Nested regression for independent value; EXP-019)};
    \node [block, below of=nest, node distance=1.8cm] (fail) {\textbf{Layer 6: Physical Failure Taxonomy} \\ \small (Isolating giant star noise \& cadence limits)};
    
    \path [line] (acc) -- (cal);
    \path [line] (cal) -- (rob);
    \path [line] (rob) -- (dom);
    \path [line] (dom) -- (nest);
    \path [line] (nest) -- (fail);
    
\end{tikzpicture}
\caption{The Multi-Dimensional Validation Hierarchy. Exoplanet vetting systems must be certified across all six layers of this protocol to ensure scientific reliability, going beyond standard single-split AUROC testing. Each certification layer maps directly to our 13 coordinated validation experiments (Section~\ref{sec:method_experiment_matrix}, Table~\ref{tab:experiment_master_matrix}): Layer 1 (Discrimination \& Accuracy): EXP-001, EXP-004, EXP-005, EXP-011, EXP-012; Layer 2 (Probability Calibration): EXP-007, EXP-009, EXP-011; Layer 3 (Measurement Robustness): EXP-007, EXP-023; Layer 4 (Domain Generalization): EXP-013; Layer 5 (Nested Control Audits): EXP-019; Layer 6 (Physical Failure Taxonomy): EXP-002, EXP-003, and host-star cohort partitioning.}
\label{fig:validation_hierarchy}
\end{figure*}

Under this protocol, we evaluate the Transit Ambiguity Recovery System (TARS) as a baseline case study. TARS employs a graph-theoretic Recovery Ambiguity Index (RAI) constructed as a parsimonious, signed sum of five standardized graph-theoretic and harmonic features derived from information theory \citep{Shannon1948, Cover2006}. To execute this multi-layer audit systematically, our study operationalizes a coordinated suite of 13 executed validation experiments selected from our 23-stage preregistered experimental registry (spanning designations EXP-001 through EXP-023, with intermediate numbers corresponding to internal engineering calibrations and exploratory runs archived during pipeline development; fully cataloged in Section~\ref{sec:method_experiment_matrix}, Table~\ref{tab:experiment_master_matrix}), spanning full-scale production runs, synthetic injection sweeps, nested feature audits, and cross-mission stress tests. We deploy this validation protocol to evaluate TARS on unblinded cohorts from TESS, Kepler, and K2 ($N=1780$, $100$, and $41$ respectively), showing that while TARS achieves comparable discrimination to complex stacks on standard splits, our multi-dimensional auditing successfully diagnoses its operational limits:
\begin{itemize}
    \item \textbf{Out-of-Domain Collapse}: Evaluating on Kepler long-cadence data reveals a collapse in ranking performance ($\text{AUROC} \approx 0.48$) and severe calibration decay (ECE $= 0.5212$), tracing a physical sensitivity limit induced by cadence-driven transit-edge blurring.
    \item \textbf{Nested Controls}: Nested regression controls (EXP-019; Section~\ref{sec:results_exp019}) expose that the graph-theoretic RAI lacks statistically significant independent predictive value once controlling for BLS SDE, directly identifying a feature redundancy that conventional evaluations fail to detect.
    \item \textbf{Cohort Boundaries}: Subgroup partitioning isolates stellar convective noise on giant hosts ($\log g < 4.0$) as a physical boundary deforming graph topology, defining where the classifier fails.
\end{itemize}

By proposing and testing this protocol, we demonstrate that rigorous evaluation can systematically reveal hidden limitations in exoplanet vetting pipelines, providing a reusable framework for certifying future automated search engines before large-scale production deployment.
